\section{Traces on varying planar continua}\label{sec:trace-analysis}

The cell construction has identified the interior equation for a limiting
mean response. To identify that response in a domain cut by observed sign
strings, we must also recover the zero boundary values on their limiting
closed set. This set may have positive area, so we must identify the
response almost everywhere on this set as well as in its open complement. The following
analytic argument uses Sobolev bounds, weighted errors at observed sites,
and the geometry of connected observed sets to identify the Dirichlet
Poisson solution.

Fix a bounded square $Q$ and a larger square $Q^*$ with
$\overline Q\subset Q^*$.  All Sobolev functions below are defined on $Q^*$;
only their restrictions to $Q$ are relevant.  This convention removes
extension issues when taking averages near $\partial Q$.  Choose a
nonnegative smooth radial function $\varphi$ supported in the unit disk,
of integral one, and write $P_ru=\varphi_r*u$ wherever the averaging disk
lies in $Q^*$.  The Dirichlet Green operator on an arbitrary open set
$V\subset Q$ is denoted by $G_V=(-\Delta_V)^{-1}$.  Its solutions and its
quadratic form are extended by zero outside $V$.

\begin{definition}[Thickness of an observed set]\label{t:thickness}
A compact set $K\subset\overline Q$ containing $\partial Q$ has
\emph{continuum thickness below $r_0$} if, for each $x\in K$ and
$0<r<r_0$, there is a compact connected set
\[
 x\in L\subset K\cap\overline B(x,2r),\qquad
 \operatorname{diam}L\ge r.
\]
For mesh approximants the same condition is required only for
$c/n\le r<r_0$, with fixed $c$.  All uniform assertions about a family of
observed sets use common constants $c,r_0$.
\end{definition}

A connected compact set of diameter at least $4r_0$ has this property:
the component through $x$ of its intersection with a closed ball meets the
boundary of that ball whenever the original continuum leaves the ball.
This is the elementary boundary-bumping property of continua.  Thus a
fixed finite union of continua whose diameters have a common positive
lower bound also qualifies, after decreasing $r_0$.  Connectedness of the
whole observed set $K$ is sufficient after decreasing $r_0$, but not
necessary.  In a localized
application one works below the lower bound on the diameters of the
components that meet the window.  A singleton observation does not qualify.

\begin{definition}[Frostman testing measures]\label{t:frostman-definition}
For $0<\alpha\le1$ and $M<\infty$, let $\mathcal M_\alpha(M)$ consist of
positive measures supported in $\overline Q$ satisfying
\[
 \mu(\overline Q)\le M,\qquad
 \mu(B(x,r))\le Mr^\alpha\quad(0<r\le1).
\]
The statement that a Sobolev function has zero $\alpha$-Frostman trace on
$K$ means that its trace vanishes in $L^1(\mu)$ for every
$\mu\in\bigcup_{M<\infty}\mathcal M_\alpha(M)$ supported on $K$.
\end{definition}

\begin{lemma}[Quantitative averaged traces]\label{t:frostman}
Let $1<p<2$, $0<\alpha\le1$, and $p+\alpha>2$.  Every
$u\in W^{1,p}(Q^*)$ has a trace $T_\mu u\in L^1(\mu)$ for
$\mu\in\mathcal M_\alpha(M)$, and
\begin{equation}\label{t:trace-rate}
 \int|T_\mu u-P_ru|\dd\mu
 \le C r^\beta\|\nabla u\|_{L^p(Q^*)},\qquad
 \beta=1-\frac{2-\alpha}{p}>0.
\end{equation}
The constant is uniform over $\mathcal M_\alpha(M)$.

For a piecewise affine mesh function $u_n$, suppose instead that a measure
on mesh vertices satisfies
\begin{equation}\label{t:coarse-frostman}
 \mu_n(B(x,r))\le C(r+\ell_n/n)^\alpha,
 \qquad \mu_n(\overline Q)\le C,
\end{equation}
where $1\le\ell_n\le C(\log n)^C$.  Then, for fixed $r>0$,
\begin{equation}\label{t:mesh-trace-rate}
 \int|u_n-P_ru_n|\dd\mu_n
 \le C\left(r^\beta+
       \ell_n^{\alpha/p}n^{-\beta}\right)
       \|\nabla u_n\|_{L^p(Q^*)}.
\end{equation}
\end{lemma}
\begin{proof}
Poincare's inequality for two nested averaging disks gives
\[
 |P_su(x)-P_{s/2}u(x)|
 \le Cs\fint_{B(x,Cs)}|\nabla u(y)|\dd y.
\]
Holder's inequality with respect to $\mu$, followed by Fubini, bounds its
integral by
\[
 Cs\left(\int\fint_{B(x,Cs)}|\nabla u(y)|^p\dd y\dd\mu(x)\right)^{1/p}
 \le Cs^{1+(\alpha-2)/p}\|\nabla u\|_p.
\]
The positive exponent makes the dyadic sum converge.  Its limit constructs
the trace and proves \eqref{t:trace-rate}; smooth approximation shows that
the construction is independent of the chosen averaging kernel.

For \eqref{t:mesh-trace-rate}, use the same sum down to the mesh scale.
Above $\ell_n/n$ it is bounded as before.  Below that scale its terms are
at most
$Cs^{1-2/p}(\ell_n/n)^{\alpha/p}\|\nabla u_n\|_p$.
Their sum down to $1/n$ is bounded by the smallest-scale term, namely
$C\ell_n^{\alpha/p}n^{-\beta}\|\nabla u_n\|_p$.
The difference between a vertex value and its mesh-scale average has the
same bound, by affine interpolation on the bounded number of incident
triangles and the same Holder calculation.
\end{proof}

\begin{lemma}[Transporting testing measures]\label{t:transport}
Suppose $K$ has continuum thickness below $r_0$.  If
$\mu\in\mathcal M_\alpha(M)$ is supported within distance $\eta$ of $K$,
where $0<\eta<r_0/10$, there is a measure
$\nu\in\mathcal M_\alpha(CM)$ supported on $K$ and a coupling $\pi$ of
$\mu,\nu$ supported on $|x-y|\le C\eta$.  Consequently
\begin{equation}\label{t:transport-trace}
 \int|T_\mu u(x)-T_\nu u(y)|\dd\pi(x,y)
 \le C\eta^\beta\|\nabla u\|_p
\end{equation}
under the exponent assumptions of Lemma~\ref{t:frostman}.

If mesh sets $K_n$ have the common thickness constants and converge to $K$
in Hausdorff distance, every such measure supported on $K$ is the weak
limit of continuous measures in $\mathcal M_\alpha(CM)$ supported on
$K_n$.  For geometric mesh graphs, these may additionally be projected
onto their vertices, giving the common bound $C(r+n^{-1})^\alpha$.
\end{lemma}
\begin{proof}
Partition the plane into squares of side $\eta$.  For each square with
positive mass choose a nearby point of $K$ and, by thickness, a continuum
of diameter at least $\eta$ within distance $C\eta$ of the square.
Such a continuum supports a probability measure $\nu_S$ with
\[
 \nu_S(B(x,r))\le C\min\{1,r/\eta\}.
\]
To construct it, project onto a coordinate axis on which the diameter is
at least $\eta/\sqrt2$ and lift normalized Lebesgue measure on the
projection interval by choosing the point with smallest remaining
coordinate in each fiber.  Compactness makes this selection Borel.
The projection of a disk of radius $r$ has length at most $2r$.

Move the mass of each square to this measure.  If $r<\eta$, at most a
bounded number of the chosen continua can meet a specified $r$-ball,
since each is within $C\eta$ of its assigned square.  The resulting mass
is at most $CM\eta^\alpha(r/\eta)\le CM r^\alpha$.
For $r\ge\eta$, use the original mass bound in the $C\eta$ enlargement
of the ball.  This proves the measure assertion and the displacement
bound.  Compare each trace value with its $\eta$ average using
\eqref{t:trace-rate}.  For paired points the difference of their averages
is at most $C\eta\fint_{B(x,C\eta)}|\nabla u|$.  Integrating as in the
preceding proof gives \eqref{t:transport-trace}.

For the final assertion take $\eta_n\downarrow0$ larger than both the
Hausdorff distance and a fixed multiple of $n^{-1}$.  Perform the same
construction on $K_n$ at scale $\eta_n$, and project its measures onto
vertices at distance $O(n^{-1})$ only if desired.  Before this projection
the measures are genuinely Frostman at every scale.  Projection changes
the mass bound to
$C(r+n^{-1})^\alpha$ and does not change the weak limit.
\end{proof}

The next calculation records exactly how weighted common-wall errors
enter the trace.  It requires more than an estimate for free-edge
gradients.

\begin{lemma}[Weighted anchor errors]\label{t:anchors}
Let $q>1$, and let $S_n$ be retained observed vertices with positive
weights $w_x$.  Suppose
\[
 \sum_{x\in S_n}w_x|z_x|^2\le E_n,
 \qquad n^{-2}\sum_{x\in S_n}w_x^{-q}\le C_q.
\]
For a probability measure $\mu_n$ on these vertices with
$a_n=\max_x\mu_n(\{x\})$,
\begin{equation}\label{t:anchor-rate}
 \sum_x\mu_n(\{x\})|z_x|
 \le C_q E_n^{1/2}n^{1/q}a_n^{(1+1/q)/2}.
\end{equation}
In particular, if $a_n\le C(\ell_n/n)^\alpha$, the bound tends to zero
in expectation when $\sup_n\E E_n<\infty$ and
$q>(2-\alpha)/\alpha$.  The deterministic inverse-weight bound may hold
outside exceptional events, provided their contributions to the displayed
quantities tend to zero.
\end{lemma}
\begin{proof}
Cauchy--Schwarz gives the upper bound
$E_n^{1/2}(\sum_x\mu_n(x)^2/w_x)^{1/2}$.
Holder and $\sum_x\mu_n(x)=1$ give
\[
 \sum_x\frac{\mu_n(x)^2}{w_x}
 \le \left(\sum_xw_x^{-q}\right)^{1/q}
      \left(\sum_x\mu_n(x)^{2q/(q-1)}\right)^{(q-1)/q}
 \le C_qn^{2/q}a_n^{1+1/q}.
\]
The resulting power of $n$ is
$-\alpha/2+(2-\alpha)/(2q)<0$.  Polylogarithmic factors do not affect
convergence, and $\E E_n^{1/2}\le(\E E_n)^{1/2}$.
\end{proof}

\begin{proposition}[Varying observed sets and filled values]
\label{t:varying-trace}
Let $K_n$ have the mesh thickness property with common constants and
converge in Hausdorff distance to $K$.  Let their mesh vertices carry
random values $z_n$, and let $U_n$ be random piecewise affine extensions
on $Q^*$.  Fix $3/2<p<2$.  Assume the following bounds, uniformly in $n$:
\begin{enumerate}
\item $\E\|U_n\|_{W^{1,p}}\le C$, and the means $\E U_n$ have a
      subsequence converging strongly in $L^p$ to $u$;
\item every vertex on $K_n$ can be moved to a retained vertex on $K_n$
      within distance $\ell_n/n$, where
      $\ell_n\le C(\log n)^C$;
\item $U_n=z_n$ at retained vertices, and their values satisfy
      Lemma~\ref{t:anchors} for some $q>3$, with $\E E_n\le C$;
\item exceptional events in these assertions have vanishing contribution
      to the Sobolev, anchor, and fixed-scale averaged quantities used
      here.
\end{enumerate}
Then $u$ has zero $1/2$-Frostman trace on $K$ and is zero Lebesgue-almost
everywhere on $K$.  If the original interpolated values also satisfy
\begin{equation}\label{t:original-filled}
 \E\|z_n-U_n\|_{L^1(Q)}\longrightarrow0,
\end{equation}
then the original mean responses have the same distributional limit,
including for smooth tests whose supports meet $K$.
\end{proposition}
\begin{proof}
Fix a $1/2$-Frostman measure $\mu$ on $K$ and approximate it by the
deterministic measures of Lemma~\ref{t:transport}.  Move their mass to the
retained vertices by the map in the hypotheses.  The moved measures may
depend on the sample; they obey \eqref{t:coarse-frostman} with the same
deterministic constants.  Lemma~\ref{t:anchors} makes their integrals of
$|U_n|$ tend to zero in expectation.  Lemma~\ref{t:frostman} bounds the
error on replacing $U_n$ by $P_rU_n$ by $Cr^\beta+o(1)$.

Keep $r>0$ fixed before removing the random movement of the testing sites.
Its additional error is at most
\[
 C\ell_n n^{-1}\|\nabla P_rU_n\|_\infty
 \le C_r\ell_n n^{-1}\|U_n\|_{W^{1,p}},
\]
whose expectation tends to zero.  We may now take expectation with the
original deterministic testing measure.  This order is essential:
expectations of field--measure products have not been factorized.
Strong $L^p$ convergence makes $P_r\E U_n$ converge uniformly to $P_ru$.
Weak convergence of the deterministic measures therefore gives
\[
 \int|P_ru|\dd\mu\le Cr^\beta.
\]
Another application of \eqref{t:trace-rate}, followed by $r\downarrow0$,
proves the zero trace.  Lebesgue measure restricted to $K$ belongs, up to
a constant, to $\mathcal M_{1/2}(M)$ on the bounded square.  Its Sobolev
trace agrees almost everywhere with $u$, which proves the area assertion.
Finally \eqref{t:original-filled} proves the equality of the two
distributional limits.
\end{proof}

\begin{lemma}[Uniform Brownian approach measures]\label{t:exit}
Let $K$ have continuum thickness, put $V=Q\setminus K$, and fix $z\in V$.
For sufficiently small $\eps>0$, let $V_\eps$ be the component containing
$z$ of $\{x\in V:\operatorname{dist}(x,V^c)>\eps\}$.  The exit law
$\mu_{z,\eps}$ of Brownian motion from $V_\eps$ satisfies
\begin{equation}\label{t:exit-frostman}
 \mu_{z,\eps}(B(x,r))\le C_z r^{1/2},
\end{equation}
uniformly in $\eps$.  The constant is uniform for $z$ a positive distance
from $K$, and depends on the common thickness constants and $Q$.
\end{lemma}
\begin{proof}
We use the planar Beurling estimate in its equivalent inward form: if an
obstacle contains a continuum crossing the annulus between radii $r$ and
$R$, the probability, from outside the outer circle, of reaching the
inner circle before the obstacle is at most $C(r/R)^{1/2}$.
The inward form follows from the usual outward estimate by inversion;
this is the Brownian estimate used in the boundary sampling arguments of
\cite{SSdiscrete}.

Each component of the closed $\eps$ neighborhood of $K$ contains a
continuum of diameter bounded below by the thickness scale.  The same is
true of the exterior obstacle.  If an $r$-ball meets the support of the
exit law, its obstacle component therefore crosses an annulus from
radius $O(r)$ to a fixed radius
$R<c\min\{r_0,\operatorname{dist}(z,K)\}$, unless $r$ is already
comparable to $R$.  Apply the inward estimate; in the latter case total
mass one suffices.  The argument is unchanged when $V_\eps$ has several
boundary components.
\end{proof}

\begin{theorem}[Poisson uniqueness from Frostman traces]\label{t:trace}
Let $K\subset\overline Q$ contain $\partial Q$ and have continuum
thickness below $r_0$.  Let $3/2<p<2$, $A>0$, and $b\in L^\infty(Q)$.
Suppose $u\in W^{1,p}(Q^*)$ has zero $1/2$-Frostman trace on $K$ and
\[
 -A\Delta u=b\quad\hbox{in }V=Q\setminus K
\]
in the sense of distributions.  Then, on $Q$,
\begin{equation}\label{t:poisson-identification}
 u=A^{-1}G_Vb,
\end{equation}
where the right side is extended by zero on $K$.  In particular the
restriction of $u$ to $V$ belongs to $H^1_0(V)$.
\end{theorem}
\begin{proof}
Let $v=A^{-1}G_Vb$, defined variationally in $H^1_0(V)$ and extended by
zero.  The Poincare inequality inherited from $Q$ gives existence.
Its zero extension is in $W^{1,p}(Q^*)$.  Approximating it in $H^1$ by
functions in $C_c^\infty(V)$, and using Lemma~\ref{t:frostman}, proves
that $v$ has zero $1/2$-Frostman trace on $K$.

Thus $w=u-v$ is harmonic on $V$ and has the same zero traces.  Fix $z\in V$
and use Lemma~\ref{t:exit}.  Transport $\mu_{z,\eps}$ to a measure on
$K$ by Lemma~\ref{t:transport}; its support is within distance $\eps$ of
$K$.  Since $\beta=1-3/(2p)>0$, \eqref{t:transport-trace} yields
\[
 \E_z|w(B_{\tau_{V_\eps}})|
 \le C_z\eps^{1-3/(2p)}\|\nabla w\|_p\longrightarrow0.
\]
The Sobolev trace on this interior boundary equals the smooth harmonic
representative of $w$.  For each fixed $\eps$, that representative is
bounded on the compact closure of $V_\eps$.  Optional stopping, or the
mean-value characterization of harmonic measure, gives
$w(z)=\E_z w(B_{\tau_{V_\eps}})=0$ in the limit.  This works in each
component of $V$.  The zero trace tested against Lebesgue measure on $K$
also gives $u=0$ almost everywhere there, completing
\eqref{t:poisson-identification} on all of $Q$.
\end{proof}

\begin{remark}[The boundary notion and the order of limits]
The proof does not promote a $W^{1,p}$ trace directly to an $H^1$ capacity
trace.  It first proves equality with the variational Green solution; the
$H^1$ assertion follows afterwards.  In a lattice application, fix
$\eps$ and use radial ball averages of radius less than $\eps/2$ before
passing to a limit.  Their limits equal the harmonic function at their
centers by the mean-value property.  No convergence of unaveraged lattice
point values is required.

Thickness cannot simply be dropped.  On the punctured unit disk,
$\log(1/|x|)$ is harmonic, has zero outer trace, and belongs to
$W^{1,p}$ for every $p<2$.  The singleton has zero $p$-capacity, and its
neighborhood hitting probabilities decay only logarithmically, so
\eqref{t:exit-frostman} fails.  Conversely, the usual slit-tip singularity
$r^{-1/2}\sin(\theta/2)$ has gradient in $L^p$ only for $p<4/3$ and
does not satisfy the hypotheses above.
\end{remark}

